% ===========================================================================
\subsection{\file{new-design-reference/css/site.css}}
% ===========================================================================

\textbf{CSS, 575 lines.} Every visual decision: colours, typefaces, spacing,
layout, movement, and how all of it rearranges on a narrow screen. It is the
longest file here and, in the judgement of this document, the most valuable thing
in the repository.

\begin{provenance}{new-design-reference/css/site.css}
Hand-written, and it states its own lineage in its second line: \code{Personal
brand site. Shares design tokens with ai.digitalise.agency.} That is a direct
statement by the author that the naming scheme came from the sibling project, and
comparing the two token blocks confirms it: identical names, identical structure,
different values.

It arrived in the first commit \code{1f7ed12} and was substantially edited once, by
\code{05210af} (``Differentiate to orange brand''). No generator banner, no vendor
prefix block of the sort a preprocessor emits, and no minified section: it is
written to be read. \textbf{Confidence: certain for hand-written, certain for the
token lineage, which rests on the author's own comment and is corroborated by the
files.}
\end{provenance}

\subsubsection*{Why it exists}

Without this file the page contains every word it needs and looks like a 1994
document: black serif text on white, one column, no colour. Nothing would be
broken and nothing would be persuasive. The two-column hero would collapse into a
stack, and the three floating badges would appear as three small blocks of text
under the placeholder.

\subsubsection*{Line by line}

\paragraph{Lines 1 to 25: the header, and the token block.}

\begin{lstlisting}[firstnumber=1,caption={\file{css/site.css}, lines 1--25}]
/* salmanadnan.com · site.css
   Personal brand site. Shares design tokens with ai.digitalise.agency. */

/* ─── TOKENS ─── */
:root {
  --s0:    #05070d;
  --s1:    #0e1118;
  --s2:    #161b26;
  --ink1:  #f4f7fc;
  --ink2:  #c3ccdd;
  --ink3:  #93a0b8;
  --accent:  #ff6719;   /* orange, personal brand, distinct from agency blue */
  --accentH: #ff8c4a;
  --accent2: #1967ff;
  --green:   #34d399;
  --border:  #24201e;
  --borderS: #1c1814;
  --font-d: 'Space Grotesk', ui-sans-serif, system-ui, sans-serif;
  --font-b: 'Inter', ui-sans-serif, system-ui, sans-serif;
  --r:   16px;
  --rsm: 10px;
  --ease: cubic-bezier(.22,1,.36,1);
  --maxw: 1240px;
  --shadow: 0 8px 40px rgba(0,0,0,.45);
}
\end{lstlisting}

Lines 1 and 2 are a comment. In CSS, anything between \code{/*} and \code{*/} is
ignored, and unlike HTML there is only this one comment form. These two lines carry
the file's entire provenance.

Line 4 is the first of twelve section banners drawn with box-drawing characters,
the author's navigation aid.

Line 5 opens the token block. \code{:root} matches the document's outermost
element, \code{<html>}. Custom properties defined there are inherited by every
element on the page, which is exactly why token blocks are written against
\code{:root}: it is the one place a definition reaches everything.

The nineteen tokens fall into six groups:

\begin{table}[H]\centering\small
\begin{tabular}{@{}llp{6.2cm}@{}}
\toprule
Tokens & Line & What the group is for \\
\midrule
\code{--s0} to \code{--s2}      & 6--8   & Surfaces, darkest to lightest. \code{s0} is the page, \code{s1} is a card, \code{s2} would be a detail inside a card \\
\code{--ink1} to \code{--ink3}  & 9--11  & Text, brightest to dimmest: a heading, body text, a caption \\
\code{--accent}, \code{--accentH}, \code{--accent2} & 12--14 & The orange, its hover variant, and the agency blue kept as a secondary \\
\code{--green}                  & 15     & The one success colour \\
\code{--border}, \code{--borderS} & 16--17 & A visible hairline, and a subtler one \\
\code{--font-d}, \code{--font-b} & 18--19 & Display typeface and body typeface \\
\code{--r}, \code{--rsm}, \code{--ease}, \code{--maxw}, \code{--shadow} & 20--24 & Corner radius, small radius, motion curve, content width, drop shadow \\
\bottomrule
\end{tabular}
\caption{The token vocabulary, identical in naming to the sibling repository and
different in almost every value.}
\end{table}

\begin{keyidea}{Why a numbered scale beats named colours}
Calling a colour \code{--s1} instead of \code{--dark-navy} names what it is
\emph{for} (the surface one step up from the page) rather than how it looks. This
pair of repositories proves the point better than any argument could: the two files
share every token \emph{name} and differ in nearly every token \emph{value}, so the
same stylesheet vocabulary produced a blue agency site and an orange personal one
with no renaming at all.
\end{keyidea}

Line 12 carries the rebrand, and it carries its own explanation in an inline
comment: \code{orange, personal brand, distinct from agency blue}. Line 14 then
keeps the agency blue as \code{--accent2}.

\begin{honest}{Four tokens are defined and never read}
\code{--s2} (line 8), \code{--accent2} (line 14), \code{--rsm} (line 21) and
\code{--shadow} (line 24) are read zero times in this file. Counted mechanically:
\code{grep -c 'var(--accent2)'} returns \code{0}, and likewise for the other three.

\code{--accent2} is the most interesting of the four. It holds the agency blue,
deliberately kept as a secondary colour, and is then never used as one. Meanwhile
five other rules in this file hard-code that same blue as a literal rather than
reading this token. The colour is used and the token for it is not, which is the
worst of both worlds and is the subject of section 7.2's largest finding, below.
\end{honest}

Line 22 defines the motion curve. \code{cubic-bezier(.22,1,.36,1)} describes how an
animation's speed varies across its run, by giving two control points of a curve.
This one starts fast and settles gently, the shape usually described as ``ease
out''. Every transition in the file refers to it, so all movement shares one
personality.

Line 23 sets the content width to 1240 pixels, 80 narrower than the sibling's 1320,
which is the kind of small deliberate difference that shows these files were tuned
separately rather than copied wholesale.

\paragraph{Lines 26 to 43: the reset, and the one utility class.}

\begin{lstlisting}[firstnumber=26,caption={\file{css/site.css}, lines 26--43}]

*, *::before, *::after { box-sizing: border-box; margin: 0; padding: 0; }
html { color-scheme: dark; scroll-behavior: smooth; }
body {
  font-family: var(--font-b);
  background: var(--s0);
  color: var(--ink1);
  -webkit-font-smoothing: antialiased;
  text-rendering: optimizeLegibility;
  overflow-x: hidden;
  line-height: 1.5;
}
img, video, svg { max-width: 100%; display: block; }
a { color: inherit; text-decoration: none; }
button { cursor: pointer; border: none; font-family: inherit; }

/* ─── UTILITY ─── */
.container { max-width: var(--maxw); margin: 0 auto; padding: 0 28px; }
\end{lstlisting}

Line 27 is a reset doing three jobs. The selector \code{*} means every element;
\code{*::before} and \code{*::after} extend that to the two invisible
pseudo-elements CSS lets you attach to anything.

\code{box-sizing: border-box} is the most consequential declaration in the file. By
default, an element given a width of 200 pixels and 20 pixels of padding occupies
240: the padding is added outside the stated width. That default surprises
everyone. \code{border-box} makes the stated width final, with padding eating into
it. Every layout calculation here assumes it.

\code{margin: 0; padding: 0} removes the browser's built-in spacing on headings and
paragraphs, so all spacing in the design is deliberate rather than inherited from
defaults that vary between browsers.

Line 28 does two things to the document root. \code{color-scheme: dark} tells the
browser this page is dark, so scrollbars, selection highlights and form controls
are drawn in their dark variants rather than staying stubbornly white.
\code{scroll-behavior: smooth} makes every jump-link glide rather than snap.

\begin{honest}{Smooth scrolling is implemented twice}
Line 28 already makes anchor links glide. \file{main.js} lines 55 to 62 then
intercept every anchor link and call \code{scrollIntoView(\{behavior: 'smooth'\})}
to do the same again. Either alone suffices. The JavaScript copy is the redundant
one and it also carries a latent bug, described in section 7.3.
\end{honest}

Lines 29 to 37 set the page defaults: body typeface, darkest surface as background,
brightest ink as text, and a line height of 1.5, meaning each line occupies one and
a half times the font size. That last number is the most important typographic
decision in the file and the one most often set badly.

\code{-webkit-font-smoothing: antialiased} on line 33 thins text rendering on Mac
displays. It is a Safari and Chrome specific property with no standard equivalent,
and it is the conventional companion to light text on a dark background, which
otherwise looks heavier than intended.

\code{overflow-x: hidden} on line 35 forbids horizontal scrolling, and here it is
load-bearing rather than cosmetic: the hero is a two-column grid whose left column
is a full-height photo frame, and the floating badges are positioned by percentage
inside it.

Line 38 makes images, videos and SVGs never exceed their container and makes them
block-level. That second part fixes a genuinely confusing default: images are
inline by default, which reserves space beneath them for the descenders of text
that is not there, producing a mysterious few pixels of gap under every image.

Line 39 strips links of their blue colour and underline, making them inherit the
surrounding colour. That is a design choice and an accessibility trade: colour
alone now distinguishes a link from text, which fails the usual guidance. Every
link here is either a button or a footer item, so the practical harm is small.

Line 40 styles buttons, and is dead in the strict sense that this page contains no
\code{<button>} element at all. It costs nothing and would come alive the moment one
was added.

Line 43 is the only utility class: a centred column no wider than 1240 pixels with
28 pixels of breathing room each side. \code{margin: 0 auto} is the standard
centring idiom, meaning no vertical margin and equal automatic margins left and
right. Every section below the hero wraps its content in this; the hero
deliberately does not, because it spans the full window.

\paragraph{Lines 44 to 72: the button system, and the first missed rebrand.}

\begin{lstlisting}[firstnumber=44,caption={\file{css/site.css}, lines 44--72}]

.btn {
  display: inline-flex;
  align-items: center;
  gap: 8px;
  padding: 13px 28px;
  border-radius: 100px;
  font-family: var(--font-d);
  font-size: 15px;
  font-weight: 600;
  letter-spacing: -0.01em;
  transition: transform 0.18s var(--ease), box-shadow 0.18s var(--ease),
              background 0.18s var(--ease);
}
.btn-primary {
  background: var(--accent);
  color: #fff;
}
.btn-primary:hover {
  background: var(--accentH);
  transform: translateY(-2px);
  box-shadow: 0 8px 24px rgba(25,103,255,.4);
}
.btn-ghost {
  background: transparent;
  color: var(--ink2);
  border: 1px solid var(--border);
}
.btn-ghost:hover { background: var(--s1); color: var(--ink1); }
\end{lstlisting}

This is the two-class pattern the HTML uses: \code{.btn} holds everything both
buttons share and \code{.btn-primary} or \code{.btn-ghost} adds only what differs.
Add a third style and you write four lines, not thirty.

\code{display: inline-flex} on line 46 makes the button a flex container sitting in
the flow of text rather than on its own line. That matters because the primary
buttons contain both a word and an icon, and flex is what lines those up.

\begin{jargon}{Flexbox}
A layout system arranging an element's children in a row or a column and
controlling how they align and share leftover space. Set \code{display: flex} on a
parent and its children become flex items. \code{align-items: center} lines them up
on their centres, which is how the icon in a button sits level with the text rather
than on the baseline.
\end{jargon}

\code{gap: 8px} on line 48 puts eight pixels between word and icon without either
needing a margin. Before \code{gap} existed this required fiddly rules targeting
all-but-the-first child, and its arrival is one of the genuine improvements in
modern CSS.

\code{border-radius: 100px} on line 50 makes a pill. Any radius larger than half the
element's height gives a perfect semicircular end, so 100 is simply comfortably
larger than needed.

Lines 55 and 56 declare the transition: three named properties, each over 0.18
seconds, each using the shared curve. Listing properties explicitly rather than
writing \code{transition: all} is better practice, because \code{all} makes the
browser watch every property and can animate things you never intended.

\begin{honest}{An orange button that casts a blue shadow}
Line 65 sets the hover shadow to \code{rgba(25,103,255,.4)}. That is
\code{\#1967ff}, the agency blue, at 40\% opacity. The button itself is orange
(\code{--accent}, line 12).

Orange and blue sit close to opposite each other on the colour wheel, so this does
not read as a subtle mismatch; it reads as a mistake. It is the most visible of the
five surviving blue literals catalogued below, and it is a one-line fix: the value
should be \code{rgba(255,103,25,.4)}, or better still derived from the token.

The cause is legible from the history. Commit \code{05210af} changed the tokens to
orange; this line never used a token, so the change could not reach it.
\end{honest}

\begin{jargon}{rgba and opacity}
\code{rgba(25,103,255,.4)} names a colour by red, green and blue amounts from 0 to
255, plus an opacity from 0 (invisible) to 1 (solid). It is the same blue as
\code{\#1967ff}, expressed in the form that allows transparency.
\end{jargon}

\begin{honest}{Neither button has a focus state}
Neither \code{.btn-primary} nor \code{.btn-ghost} defines \code{:focus} or
\code{:focus-visible}. A visitor navigating by keyboard rather than mouse gets
whatever outline the browser draws by default, which on a dark background is often
nearly invisible. Since these are the page's only interactive elements, this is the
most consequential accessibility gap in the stylesheet, and it is three lines to
fix.
\end{honest}

\paragraph{Lines 73 to 88: the two reveal mechanisms.}

\begin{lstlisting}[firstnumber=73,caption={\file{css/site.css}, lines 73--88}]

/* ─── REVEAL ─── */
.reveal {
  opacity: 0;
  transform: translateY(28px);
  transition: opacity 0.7s var(--ease), transform 0.7s var(--ease);
}
.reveal.visible { opacity: 1; transform: none; }

.stagger > * {
  opacity: 0;
  transform: translateY(20px);
  transition: opacity 0.6s var(--ease), transform 0.6s var(--ease);
  transition-delay: calc(var(--i, 0) * 90ms);
}
.stagger.visible > * { opacity: 1; transform: none; }
\end{lstlisting}

These sixteen lines are the entire animation contract between the stylesheet and
the JavaScript: the CSS defines both states and the transition between them, and
the JavaScript's only job is to add the class \code{visible} at the right moment.

\code{.reveal} starts invisible and 28 pixels low. \code{.reveal.visible}, which
matches an element carrying \emph{both} classes, returns it to fully visible and
its natural position. Because a transition is declared, the browser animates
between the two over 0.7 seconds rather than snapping.

The stagger is the same idea with one addition. \code{.stagger > *} means every
direct child of a staggered container, and line 86 delays each child by
\code{calc(var(--i, 0) * 90ms)}: the number the HTML wrote inline, multiplied by 90
milliseconds. Child 0 starts immediately, child 5 starts 450 milliseconds later.

\begin{jargon}{calc, and a custom property's fallback}
\code{calc()} performs arithmetic in CSS and can mix units. \code{var(--i, 0)}
reads the custom property \code{--i} and falls back to \code{0} if it was never
set, which keeps this rule safe on any element whose inline \code{--i} was
forgotten: it gets no delay instead of breaking.
\end{jargon}

\begin{keyidea}{Why the animation lives in CSS and the trigger lives in JavaScript}
This split is why the page degrades well, and it is also its one real hazard. If
JavaScript never runs, the class \code{visible} is never added and every revealed
element stays at opacity zero, invisible forever. That would be catastrophic, and
it is exactly why line 573 exists: the reduced-motion block at the end of the file
forces \code{.reveal} and \code{.stagger > *} back to full opacity. That block is
doing double duty as an accessibility feature and as a safety net.
\end{keyidea}

\paragraph{Lines 89 to 122: the navigation bar.}

\begin{lstlisting}[firstnumber=89,caption={\file{css/site.css}, lines 89--122}]

/* ─── NAV ─── */
.nav {
  position: fixed;
  top: 0; left: 0; right: 0;
  z-index: 100;
  padding: 18px 28px;
  display: flex;
  align-items: center;
  justify-content: space-between;
  transition: background 0.3s, border-color 0.3s, box-shadow 0.3s;
}
.nav.scrolled {
  background: rgba(5,7,13,0.88);
  border-bottom: 1px solid var(--border);
  box-shadow: 0 4px 24px rgba(0,0,0,.3);
  backdrop-filter: blur(16px);
  -webkit-backdrop-filter: blur(16px);
}
.nav-inner {
  max-width: var(--maxw);
  width: 100%;
  margin: 0 auto;
  display: flex;
  align-items: center;
  justify-content: space-between;
}
.nav-logo {
  font-family: var(--font-d);
  font-size: 17px;
  font-weight: 700;
  letter-spacing: -0.02em;
  color: var(--ink1);
}
\end{lstlisting}

\code{position: fixed} on line 92 pins the bar to the window rather than the
document, so it stays put while the page scrolls beneath it. Line 93 pins it to the
top and stretches it across.

\code{z-index: 100} on line 94 controls stacking order: when elements overlap, the
higher number is drawn in front. 100 is arbitrary and generous, chosen so nothing
else can accidentally cover the bar.

Lines 101 to 107 are the scrolled state, the class \file{main.js} adds once the page
has moved more than 40 pixels: a background at 88\% opacity, a bottom hairline, and
a shadow.

\code{backdrop-filter: blur(16px)} on line 105 is the interesting one: it blurs
whatever is \emph{behind} the element, producing the frosted-glass effect that only
works because the background is slightly transparent. Line 106 repeats it with a
\code{-webkit-} prefix for Safari, which needed the prefixed form for years.

\begin{honest}{A doubled flex layout}
\code{.nav} sets \code{display: flex} with \code{justify-content: space-between}
(lines 96 to 98), and \code{.nav-inner} sets the same three properties again (lines
112 to 114). The outer one has a single child, so its \code{space-between} has
nothing to space. It is harmless and redundant, and it makes the layout harder to
reason about than it needs to be. The sibling repository has exactly the same
doubling, so it came from the shared skeleton.
\end{honest}

Lines 116 to 122 style the logo. Note there is no rule for a logo mark or an
accent-coloured fragment, unlike the sibling: this logo is one span of plain text,
matching the simpler markup described in section 7.1.

\paragraph{Lines 123 to 133: the hero as a two-column grid.}

\begin{lstlisting}[firstnumber=123,caption={\file{css/site.css}, lines 123--133}]

/* ─── HERO ─── */
.hero {
  min-height: 100vh;
  display: grid;
  grid-template-columns: 1fr 1fr;
  align-items: center;
  gap: 0;
  padding-top: 80px;
  overflow: hidden;
}
\end{lstlisting}

This is where the two siblings genuinely diverge. The other repository's hero is a
single centred column over a field of drifting cards; this one is a two-column
grid, photograph on the left and words on the right.

\begin{jargon}{CSS Grid}
A layout system for two dimensions, as opposed to flexbox's one.
\code{grid-template-columns: 1fr 1fr} means two columns each taking one equal
share of the space. \code{1fr} is a fraction of what remains after any fixed sizes
are accounted for.
\end{jargon}

\code{min-height: 100vh} on line 126 makes the hero at least as tall as the window.
\code{vh} means one percent of the viewport height. \code{padding-top: 80px} on
line 131 leaves room for the fixed navigation bar, which is removed from the normal
flow and would otherwise sit on top of the content.

\paragraph{Lines 134 to 148: the photo column.}

\begin{lstlisting}[firstnumber=134,caption={\file{css/site.css}, lines 134--148}]

/* photo side */
.hero-photo-side {
  position: relative;
  height: 100vh;
  display: flex;
  align-items: center;
  justify-content: center;
  overflow: hidden;
}
.hero-photo-frame {
  position: relative;
  width: 100%;
  height: 100%;
}
\end{lstlisting}

\code{position: relative} on line 137 looks like it does nothing and is essential.
It establishes this element as the reference point for any absolutely positioned
descendant, which here means the three floating badges. Without it they would
position themselves against the whole window instead of against the photo column.

\begin{jargon}{Positioning: relative and absolute}
\code{position: absolute} takes an element out of the normal flow and places it at
coordinates you give, measured from the nearest ancestor that is itself positioned,
which in practice means the nearest one with \code{position: relative}. This
pairing, a relative parent and absolute children, is the standard way to place
things inside a box.
\end{jargon}

Line 138 sets the column's height to \code{100vh}, the full window height, which is
what makes the photograph bleed to the top and bottom edges of the screen.

\paragraph{Lines 149 to 179: the placeholder, styled with real care.}

\begin{lstlisting}[firstnumber=149,caption={\file{css/site.css}, lines 149--179}]

/* Photo placeholder - replace with <img src="assets/salman.jpg"> */
.hero-photo-placeholder {
  width: 100%;
  height: 100%;
  background: linear-gradient(160deg, #1a0800 0%, #4a1800 40%, #05070d 100%);
  display: flex;
  align-items: center;
  justify-content: center;
  flex-direction: column;
  gap: 16px;
  color: rgba(255,255,255,.15);
  font-family: var(--font-d);
}
.hero-photo-placeholder .placeholder-initial {
  font-size: 120px;
  font-weight: 700;
  letter-spacing: -0.05em;
  line-height: 1;
  background: linear-gradient(135deg, rgba(255,103,25,.6), rgba(255,140,74,.3));
  -webkit-background-clip: text;
  -webkit-text-fill-color: transparent;
  background-clip: text;
}
.hero-photo-placeholder .placeholder-note {
  font-size: 13px;
  color: rgba(255,255,255,.25);
  text-align: center;
  max-width: 220px;
  line-height: 1.5;
}
\end{lstlisting}

These thirty lines style the placeholder that section 7.1 identified as the most
consequential thing in the repository, and they are worth reading for what they say
about intent.

Line 154 fills the frame with a three-stop gradient running from a near-black
brown, through a deep orange, to the page background. That is a considered stand-in
for a photograph: it establishes the brand colour and fades into the page rather
than ending at a hard edge.

Lines 163 to 172 style the giant letter, and they use a genuinely advanced
technique.

\begin{keyidea}{How the letter is filled with a gradient}
You cannot pour a gradient into text directly. The trick on lines 168 to 171 is
three steps: give the element a gradient background, tell the browser to clip that
background to the shape of the text (\code{background-clip: text}), and then make
the text itself transparent (\code{-webkit-text-fill-color: transparent}) so the
clipped background shows through the letterforms.

Line 171 sets the standard \code{background-clip: text} and line 169 the prefixed
\code{-webkit-} form, which is still needed for full support. Note the third line
has no unprefixed equivalent at all: \code{-webkit-text-fill-color} is a
non-standard property that every major browser implements, and this effect depends
on it.
\end{keyidea}

Lines 173 to 179 style the instruction text: 13 pixels, white at 25\% opacity,
centred, capped at 220 pixels wide.

\begin{honest}{The placeholder note is deliberately styled to be readable}
This is the detail that settles how to describe the placeholder. The note is not
raw unstyled debris: it has been given a size, a colour, an alignment, a maximum
width and a line height. Somebody laid it out to be read.

Read charitably, that is a designer making an unfinished state look tidy in a
review. Read plainly, it means a visitor to the published page would see a legible
instruction addressed to a developer. Both readings are consistent with the
evidence and the second is what a visitor experiences, which is why section 14
treats it as the project's defining unfinished element.
\end{honest}

\paragraph{Lines 180 to 200: the two edge fades.}

\begin{lstlisting}[firstnumber=180,caption={\file{css/site.css}, lines 180--200}]

/* fade gradient on the right edge so text side blends */
.hero-photo-frame::after {
  content: '';
  position: absolute;
  top: 0; right: 0;
  width: 40%;
  height: 100%;
  background: linear-gradient(to right, transparent, var(--s0));
  pointer-events: none;
}
/* fade at bottom */
.hero-photo-frame::before {
  content: '';
  position: absolute;
  bottom: 0; left: 0; right: 0;
  height: 30%;
  background: linear-gradient(to top, var(--s0), transparent);
  pointer-events: none;
  z-index: 1;
}
\end{lstlisting}

Two pseudo-elements, and they solve the problem a hard-edged half-page image
creates.

\begin{jargon}{Pseudo-element}
An extra box CSS creates inside an element with no corresponding tag in the HTML,
written \code{::before} or \code{::after}. It needs the \code{content} property to
exist at all, which is why lines 183 and 193 set \code{content: ''}: an empty
string is enough to bring it into being. It is the standard way to draw a
decorative shape without polluting the markup.
\end{jargon}

The \code{::after} (lines 182 to 190) covers the rightmost 40\% of the photo column
with a gradient running from transparent to the page background, so the photograph
dissolves into the text column instead of ending at a visible seam. The
\code{::before} (lines 192 to 200) does the same across the bottom 30\%.

Both carry \code{pointer-events: none}, which makes them transparent to the mouse
so clicks pass through. Without it these two invisible layers would sit over the
photo column and swallow any interaction beneath them.

\begin{keyidea}{Why the bottom fade needs a z-index and the right fade does not}
Line 199 gives the \code{::before} a \code{z-index: 1}. The reason is ordering:
without an explicit stacking position, a \code{::before} is painted \emph{behind}
its element's content, while an \code{::after} is painted in front. So the right
fade works by default and the bottom fade would be hidden behind the placeholder
until lifted.

The badges then sit at \code{z-index: 2} (line 214), one layer above the bottom
fade, which is what keeps them crisp while the photograph dims underneath. Three
elements, three deliberate layers.
\end{keyidea}

\paragraph{Lines 201 to 241: the floating badges.}

\begin{lstlisting}[firstnumber=201,caption={\file{css/site.css}, lines 201--241}]

/* floating badges on the photo */
.hero-float-badge {
  position: absolute;
  background: rgba(13,18,32,0.85);
  backdrop-filter: blur(12px);
  -webkit-backdrop-filter: blur(12px);
  border: 1px solid var(--border);
  border-radius: 12px;
  padding: 12px 16px;
  display: flex;
  align-items: center;
  gap: 10px;
  z-index: 2;
  animation: badgeFloat 6s ease-in-out infinite;
}
.hero-float-badge:nth-child(2) { bottom: 25%; left: 8%;  animation-delay: 0s;   animation-duration: 7s; }
.hero-float-badge:nth-child(3) { bottom: 40%; right: 12%; animation-delay: 2s;  animation-duration: 6s; }
.hero-float-badge:nth-child(4) { top: 22%;   left: 10%;  animation-delay: 1s;  animation-duration: 8s; }

@keyframes badgeFloat {
  0%, 100% { transform: translateY(0); }
  50%       { transform: translateY(-8px); }
}
.badge-icon {
  width: 36px; height: 36px;
  border-radius: 10px;
  display: flex; align-items: center; justify-content: center;
}
.badge-icon-blue   { background: rgba(25,103,255,.15); }
.badge-icon-orange { background: rgba(255,103,25,.15); }
.badge-icon-green  { background: rgba(52,211,153,.12); }
.badge-icon svg { width: 18px; height: 18px; }
.badge-text-main {
  font-family: var(--font-d);
  font-size: 14px;
  font-weight: 700;
  color: var(--ink1);
  white-space: nowrap;
}
.badge-text-sub { font-size: 11px; color: var(--ink3); white-space: nowrap; }
\end{lstlisting}

Lines 203 to 216 define a badge: absolutely positioned, dark and semi-transparent
with a 12-pixel blur behind it, bordered, rounded, laid out as a row with a gap,
lifted to \code{z-index: 2}, and animated forever.

Note line 205: \code{rgba(13,18,32,0.85)}.

\begin{honest}{A colour left behind by the rebrand, in a subtler way}
\code{rgb(13,18,32)} is \code{\#0d1220}, which is the \emph{sibling repository's}
\code{--s1} surface colour. This file's own \code{--s1} is \code{\#0e1118}, a
slightly different, less blue dark. So the badge background is the other project's
card colour, written as a literal, at 85\% opacity.

Nobody would notice: the two values differ by a few points of blue and the badge
sits over a gradient at 85\% opacity. It is included here because it is evidence
about how the rebrand was done, not because it is a visible defect. The tokens were
updated and every literal in the file was left exactly as inherited.
\end{honest}

Lines 217 to 219 place the three badges individually, and this is where the HTML's
child order becomes load-bearing. \code{:nth-child(2)} matches the second child of
its parent, so these three rules assume the placeholder occupies child 1, as section
7.1 explained. Each badge gets a different corner, a different animation delay and a
different duration (7, 6 and 8 seconds), which is what stops the three drifting in
visible unison.

Lines 221 to 224 define the drift itself: a gentle rise of 8 pixels at the halfway
point and back. Grouping \code{0\%, 100\%} on one line is the idiom for an
animation returning to where it started.

\begin{jargon}{Keyframes}
An animation defined by describing the element at particular points through its
run, given as percentages of the total duration. The browser fills in every moment
between. \code{0\%} is the start and \code{100\%} the end.
\end{jargon}

Lines 230 to 232 are the three tinted icon chips: blue, orange and green, each the
matching colour at 12 to 15\% opacity. As section 7.1 noted, these three are the one
place blue is plainly intentional.

Line 239 and line 241 both set \code{white-space: nowrap}, stopping ``1.6M+ views''
and ``organic, for clients'' wrapping onto two lines and making the badge grow
taller than its neighbours.

\paragraph{Lines 242 to 301: the content column.}

\begin{lstlisting}[firstnumber=242,caption={\file{css/site.css}, lines 242--301}]

/* content side */
.hero-content-side {
  padding: 0 48px 0 20px;
  display: flex;
  flex-direction: column;
  justify-content: center;
  gap: 32px;
}
.hero-tag {
  display: inline-flex;
  align-items: center;
  gap: 8px;
  background: rgba(25,103,255,.1);
  border: 1px solid rgba(25,103,255,.25);
  border-radius: 100px;
  padding: 6px 14px;
  font-size: 13px;
  font-weight: 600;
  color: var(--accentH);
  width: fit-content;
}
.hero-h1 {
  font-family: var(--font-d);
  font-size: clamp(40px, 5vw, 70px);
  font-weight: 700;
  letter-spacing: -0.03em;
  line-height: 1.04;
  color: var(--ink1);
}
.hero-h1 em { font-style: normal; color: var(--accent); }
.hero-descriptor {
  font-family: var(--font-d);
  font-size: 20px;
  font-weight: 600;
  color: var(--ink3);
  letter-spacing: -0.01em;
}
.hero-actions {
  display: flex;
  align-items: center;
  gap: 12px;
  flex-wrap: wrap;
}
.hero-mini-stats {
  display: flex;
  align-items: center;
  gap: 28px;
  padding-top: 4px;
  border-top: 1px solid var(--borderS);
}
.hero-mini-stat { text-align: left; }
.hero-mini-num {
  font-family: var(--font-d);
  font-size: 24px;
  font-weight: 700;
  color: var(--ink1);
  letter-spacing: -0.03em;
}
.hero-mini-lbl { font-size: 12px; color: var(--ink3); }
\end{lstlisting}

Line 244 lays the right column out as a vertical flex column with a 32-pixel gap,
so its five children space themselves without margins. The asymmetric padding
(\code{0 48px 0 20px}) pulls the text slightly left, toward the fading edge of the
photograph.

Lines 251 to 263 are the small pill above the name, and it contains the second
missed rebrand.

\begin{honest}{A blue pill containing orange text}
Lines 255 and 256 tint the pill with \code{rgba(25,103,255,...)}, the agency blue,
while line 261 colours its text \code{var(--accentH)}, which is the orange
\code{\#ff8c4a}.

So the element reads its foreground from a token and its background from a
literal, and after the rebrand the two no longer agree. On screen it is a
blue-tinted capsule with orange text inside it, which looks like neither brand.

\code{width: fit-content} on line 262 is worth noting separately and is correct: in
a flex column, a child stretches to full width by default, so without this the pill
would span the entire column instead of hugging its text.
\end{honest}

Line 266 is the most useful single line in the file to learn from.

\begin{jargon}{clamp}
\code{clamp(minimum, preferred, maximum)} takes the preferred value but never lets
it fall below the minimum or rise above the maximum. Here the preferred value is
\code{5vw}, 5\% of the window width, so the headline scales continuously between a
floor of 40 pixels and a ceiling of 70. One line replaces what used to need several
media queries.
\end{jargon}

Line 268 sets \code{letter-spacing: -0.03em}, pulling letters closer together. Large
display type looks loose at its default spacing, and negative tracking at large
sizes is standard practice. The unit \code{em} means ``relative to this element's
font size'', so the tightening scales with the text.

Line 272 is the rule flagged in section 7.1: it takes the \code{<em>} tag, removes
the italics it normally implies, and colours it with the accent instead.

Lines 286 to 301 build the three small statistics: a flex row with a 28-pixel gap
above a hairline, each statistic a 24-pixel number over a 12-pixel label.

\paragraph{Lines 302 to 321: the work section, and the shared headings.}

\begin{lstlisting}[firstnumber=302,caption={\file{css/site.css}, lines 302--321}]

/* ─── WORK GRID ─── */
.work-section { padding: 120px 0 80px; }
.section-kicker {
  font-size: 12px;
  font-weight: 700;
  text-transform: uppercase;
  letter-spacing: 0.1em;
  color: var(--accent);
  margin-bottom: 14px;
}
.section-title {
  font-family: var(--font-d);
  font-size: clamp(28px, 4vw, 48px);
  font-weight: 700;
  letter-spacing: -0.025em;
  line-height: 1.1;
  color: var(--ink1);
  margin-bottom: 52px;
}
\end{lstlisting}

Although they sit under the work banner, \code{.section-kicker} and
\code{.section-title} are used by three sections. Their placement here is a small
organisational untidiness: a reader looking for the heading styles would reasonably
expect them near the top with the other shared pieces.

The kicker is the small uppercase line (``The work'', ``Case studies'') and the
title is the large heading beneath it. Note the tracking is positive on the kicker
and negative on the title. That is not an inconsistency but a rule of typography:
large text needs tightening and small uppercase text needs opening up or it becomes
a smudge.

\begin{honest}{The kicker's contrast is worth measuring}
\code{--accent} is now \code{\#ff6719}, a saturated orange, set at 12 pixels
uppercase on a near-black background. Small text needs more contrast than large
text, not less. This document does not have a measured contrast ratio for that
exact pairing and so does not assert a failure, but it is the first thing to check
in an accessibility audit. Orange on near-black generally fares better than the
sibling's blue on near-black, so this is the less likely of the two to fail.
\end{honest}

\paragraph{Lines 322 to 388: the project cards.}

\begin{lstlisting}[firstnumber=322,caption={\file{css/site.css}, lines 322--388}]

.work-grid {
  display: grid;
  grid-template-columns: repeat(3, 1fr);
  gap: 14px;
}
.work-card {
  background: var(--s1);
  border: 1px solid var(--border);
  border-radius: var(--r);
  overflow: hidden;
  transition: transform 0.22s var(--ease), border-color 0.22s, box-shadow 0.22s;
}
.work-card:hover {
  transform: translateY(-4px);
  border-color: rgba(25,103,255,.35);
  box-shadow: 0 16px 48px rgba(0,0,0,.4);
}
.work-thumb {
  aspect-ratio: 16/9;
  position: relative;
  display: flex; align-items: center; justify-content: center;
  overflow: hidden;
}
.work-thumb-overlay {
  position: absolute; inset: 0;
  background: rgba(5,7,13,.2);
  display: flex; align-items: center; justify-content: center;
  opacity: 0;
  transition: opacity 0.22s;
}
.work-card:hover .work-thumb-overlay { opacity: 1; }
.work-play {
  width: 48px; height: 48px;
  border-radius: 50%;
  background: rgba(255,255,255,.15);
  backdrop-filter: blur(8px);
  display: flex; align-items: center; justify-content: center;
  border: 1.5px solid rgba(255,255,255,.3);
}
.work-play svg { width: 20px; height: 20px; color: #fff; margin-left: 2px; }
.work-badge {
  position: absolute; top: 10px; right: 10px;
  background: rgba(5,7,13,.75);
  backdrop-filter: blur(8px);
  border: 1px solid var(--border);
  border-radius: 20px;
  padding: 4px 10px;
  font-size: 11px; font-weight: 600;
  color: var(--green);
}
.wt1 { background: linear-gradient(135deg,#2a1000,#5a2800); }
.wt2 { background: linear-gradient(135deg,#2a1800,#5a3000); }
.wt3 { background: linear-gradient(135deg,#1a0a00,#3a1800); }
.wt4 { background: linear-gradient(135deg,#300800,#602000); }
.wt5 { background: linear-gradient(135deg,#200e00,#482800); }
.wt6 { background: linear-gradient(135deg,#281400,#503000); }
.work-info { padding: 16px; }
.work-client {
  font-size: 11px; color: var(--ink3); font-weight: 600;
  text-transform: uppercase; letter-spacing: 0.06em; margin-bottom: 4px;
}
.work-title {
  font-family: var(--font-d);
  font-size: 15px; font-weight: 600;
  color: var(--ink1); letter-spacing: -0.01em;
}
\end{lstlisting}

\code{overflow: hidden} on line 332 is what makes the rounded corners work: without
it the gradient thumbnail inside would square off the top two corners it is
supposed to follow.

\code{aspect-ratio: 16/9} on line 341 reserves a widescreen-shaped area regardless
of width. This is worth appreciating: the layout will not shift when a real image
eventually loads, because the space is already the right shape. It is the modern
replacement for a well-known hack involving percentage padding.

The overlay (lines 346 to 353) sits over the whole thumbnail at zero opacity and
fades to one on hover, revealing the play button beneath a 20\% black wash. The play
button (lines 354 to 361) is a frosted circle: white at 15\% opacity with an
8-pixel backdrop blur and a brighter rim.

Line 362 gives the play triangle \code{margin-left: 2px}. That is not an accident.
A triangle's visual centre of mass sits left of its geometric centre, so a
perfectly centred play icon looks slightly left. Nudging it two pixels right is a
standard correction, and noticing it is a good sign that whoever built this cared
about detail.

\begin{honest}{The third missed rebrand}
Line 337 sets the hover border to \code{rgba(25,103,255,.35)}, the agency blue, on
a card in an orange-branded site. Same cause, same one-line fix as line 65.
\end{honest}

Lines 373 to 378 are the six placeholder gradients, and they \emph{were} rebranded:
every one is a two-stop orange or brown blend, where the sibling's equivalents run
through blue, teal, violet and green. So the rebrand did reach these six literals
and missed five others, which tells you it was done by eye rather than by search.

Lines 380 to 388 style the two lines of text under each thumbnail: a small uppercase
client name in the dimmest ink, and the result in the display typeface at 15
pixels. This is the only place in the file where the display typeface does ordinary
work at a small size.

\paragraph{Lines 389 to 419: the metrics band.}

\begin{lstlisting}[firstnumber=389,caption={\file{css/site.css}, lines 389--419}]

/* ─── METRICS ─── */
.metrics-section {
  padding: 80px 0;
  border-top: 1px solid var(--borderS);
  border-bottom: 1px solid var(--borderS);
}
.metrics-grid {
  display: grid;
  grid-template-columns: repeat(4, 1fr);
  gap: 2px;
}
.metric-item {
  padding: 36px 24px;
  text-align: center;
  border-right: 1px solid var(--borderS);
}
.metric-item:last-child { border-right: none; }
.metric-num {
  font-family: var(--font-d);
  font-size: clamp(34px, 4vw, 56px);
  font-weight: 700;
  letter-spacing: -0.03em;
  line-height: 1;
  color: var(--ink1);
  margin-bottom: 8px;
}
.metric-label {
  font-size: 13px; color: var(--ink3); font-weight: 500;
  text-transform: uppercase; letter-spacing: 0.06em;
}
\end{lstlisting}

Four equal columns with a 2-pixel gap, and the separators drawn as a right border
on every item (line 404) with the last one removed (line 406).

Note that this is the only place in the file that draws dividing lines this way.
The sibling repository contains a second, cleverer technique for the same job (a
grid whose gaps show a coloured layer beneath), used there on its services grid. That
technique does not appear here, because this file has no equivalent grid, so the
inconsistency the sibling has within itself does not exist in this one.

Lines 407 to 415 style the big numbers: display typeface, scaling between 34 and 56
pixels, with \code{line-height: 1}. That last is deliberate. The page-wide default
of 1.5 is right for paragraphs and wrong for a large number, where it leaves a
visible gap above and below.

\paragraph{Lines 420 to 479: the case-study cards, and three dead rules.}

\begin{lstlisting}[firstnumber=420,caption={\file{css/site.css}, lines 420--479}]

/* ─── CASE STUDIES ─── */
.cases-section { padding: 100px 0; }
.case-card {
  display: grid;
  grid-template-columns: 1fr 1fr;
  gap: 48px;
  align-items: center;
  background: var(--s1);
  border: 1px solid var(--border);
  border-radius: 24px;
  padding: 48px;
  margin-bottom: 16px;
  transition: border-color 0.22s;
}
.case-card:hover { border-color: rgba(25,103,255,.3); }
.case-card.reverse { direction: rtl; }
.case-card.reverse > * { direction: ltr; }
.case-visual {
  border-radius: var(--r);
  overflow: hidden;
  aspect-ratio: 16/10;
  position: relative;
}
.cv1 { background: linear-gradient(135deg,#2a1000,#5a2800); }
.cv2 { background: linear-gradient(135deg,#300800,#602000); }
.case-visual-inner {
  position: absolute; inset: 0;
  display: flex; align-items: flex-end;
  padding: 20px;
  background: linear-gradient(to top, rgba(5,7,13,.8) 0%, transparent 60%);
}
.case-visual-metrics {
  display: flex; gap: 12px; flex-wrap: wrap;
}
.cvm {
  background: rgba(5,7,13,.7);
  backdrop-filter: blur(8px);
  border: 1px solid var(--border);
  border-radius: 20px;
  padding: 6px 12px;
  display: flex; flex-direction: column; gap: 1px;
}
.cvm-num {
  font-family: var(--font-d);
  font-size: 16px; font-weight: 700;
  color: var(--ink1); letter-spacing: -0.02em;
}
.cvm-lbl { font-size: 10px; color: var(--ink3); text-transform: uppercase; letter-spacing: 0.06em; }
.case-content { display: flex; flex-direction: column; gap: 20px; }
.case-client-name {
  font-size: 12px; font-weight: 700;
  color: var(--accent); text-transform: uppercase; letter-spacing: 0.08em;
}
.case-title {
  font-family: var(--font-d);
  font-size: clamp(20px, 2.5vw, 30px);
  font-weight: 700; letter-spacing: -0.02em;
  color: var(--ink1); line-height: 1.2;
}
\end{lstlisting}

Each case card is a two-column grid with a 48-pixel gap: a visual panel and a block
of text, vertically centred against each other.

Lines 436 and 437 are the dead rules identified in section 7.1.

\begin{honest}{Three rules that have never applied, and the mechanism they were
replaced by}
\code{.case-card.reverse} means ``an element carrying both the class
\code{case-card} and the class \code{reverse}''. The HTML at line 321 writes
\code{class="case-card case-card-reverse reveal"}. A hyphen does not create a second
class, so this selector matches nothing, and neither does its partner on line 437 or
the third instance in the responsive block at line 559.

What they were trying to do is elegant and worth understanding even though it is
dead. Setting \code{direction: rtl} on a grid container reverses the order its
children are placed in, and setting \code{direction: ltr} on the children puts their
text back the right way round. Two lines, and a two-column layout mirrors itself.

What actually happens instead is that \file{index.html} sets \code{order:2} and
\code{order:1} inline on the two children, which achieves the same result by a
different route. So the design is correct, three rules are dead, and nothing
reported it. This is the clearest illustration in either sibling of what ``silent
failure'' means: an entire mechanism can be broken while the page looks right,
because a second mechanism quietly does the same job.
\end{honest}

\begin{honest}{The fourth missed rebrand}
Line 435 sets the case card's hover border to \code{rgba(25,103,255,.3)}, the agency
blue again.
\end{honest}

Lines 446 to 451 lay the metric chips over the bottom of each visual, with a
gradient from 80\% opaque page colour at the bottom to transparent at 60\% height, so
the chips sit on a dark base and the top of the visual stays clear. Lines 455 to 468
style the chips themselves: frosted, rounded, a number over a tiny uppercase label.

\paragraph{Lines 480 to 524: the booking panel.}

\begin{lstlisting}[firstnumber=480,caption={\file{css/site.css}, lines 480--524}]

/* ─── BOOK A CALL ─── */
.book-section { padding: 80px 0 140px; }
.book-wrap {
  background: var(--s1);
  border: 1px solid var(--border);
  border-radius: 24px;
  padding: 72px 48px;
  text-align: center;
  position: relative;
  overflow: hidden;
}
.book-wrap::before {
  content: '';
  position: absolute;
  top: -100px; left: 50%;
  transform: translateX(-50%);
  width: 600px; height: 300px;
  background: radial-gradient(ellipse, rgba(25,103,255,.1) 0%, transparent 70%);
  pointer-events: none;
}
.book-eyebrow {
  font-size: 12px; font-weight: 700;
  text-transform: uppercase; letter-spacing: 0.1em;
  color: var(--accent); margin-bottom: 16px;
}
.book-title {
  font-family: var(--font-d);
  font-size: clamp(28px, 4vw, 48px);
  font-weight: 700; letter-spacing: -0.025em;
  color: var(--ink1); margin-bottom: 12px;
}
.book-sub {
  font-size: 16px; color: var(--ink3); margin-bottom: 36px;
}
.book-cta-row {
  display: flex; align-items: center; justify-content: center;
  gap: 14px; flex-wrap: wrap;
}
.book-detail {
  display: flex; align-items: center; justify-content: center;
  gap: 6px; font-size: 13px; color: var(--ink3); margin-top: 20px;
}
.book-detail svg { width: 15px; height: 15px; }
.book-detail span { color: var(--ink2); }
\end{lstlisting}

Lines 492 to 500 add a glow above the panel using a \code{::before}: a 600 by 300
ellipse positioned 100 pixels above the panel's top edge and centred with the
standard \code{left: 50\%} plus \code{translateX(-50\%)} idiom. Because the panel has
\code{overflow: hidden}, only the lower part of the glow is visible, which makes it
read as light spilling in from above rather than as a floating oval.

\code{pointer-events: none} on line 499 appears for the third time in the file, for
the same reason as before: a decorative layer must not intercept the click aimed at
the booking button.

\begin{honest}{The fifth and last missed rebrand}
Line 498 makes that glow \code{rgba(25,103,255,.1)}, the agency blue, on the most
important element of the page. It is subtle at 10\% opacity, and it means the halo
around the call to action is the wrong brand colour.

That completes the catalogue: lines 65, 255, 256, 337, 435 and 498, five rules and
six declarations, all inherited from the sibling and all missed by a rebrand that
correctly updated the tokens and the six thumbnail gradients.
\end{honest}

\paragraph{Lines 525 to 535: the footer.}

\begin{lstlisting}[firstnumber=525,caption={\file{css/site.css}, lines 525--535}]

/* ─── FOOTER ─── */
footer {
  border-top: 1px solid var(--borderS);
  padding: 32px 28px;
  text-align: center;
  font-size: 13px;
  color: var(--ink3);
}
footer a { color: var(--ink3); transition: color 0.15s; }
footer a:hover { color: var(--ink1); }
\end{lstlisting}

Deliberately unremarkable: one hairline, centred dim text, and links that brighten
on hover over 0.15 seconds.

\paragraph{Lines 536 to 575: the responsive rules and the accessibility block.}

\begin{lstlisting}[firstnumber=536,caption={\file{css/site.css}, lines 536--575}]

/* ─── RESPONSIVE ─── */
@media (max-width: 960px) {
  .hero {
    grid-template-columns: 1fr;
    min-height: auto;
    padding-top: 80px;
  }
  .hero-photo-side {
    height: 55vw;
    min-height: 280px;
  }
  .hero-content-side {
    padding: 40px 28px 60px;
  }
  .hero-mini-stats { gap: 20px; }
  .work-grid   { grid-template-columns: repeat(2, 1fr); }
  .metrics-grid { grid-template-columns: repeat(2, 1fr); }
  .metric-item { border-right: none; border-bottom: 1px solid var(--borderS); }
  .metric-item:nth-child(odd) { border-right: 1px solid var(--borderS); }
  .metric-item:nth-child(3),
  .metric-item:nth-child(4)   { border-bottom: none; }
  .case-card { grid-template-columns: 1fr; gap: 28px; padding: 32px; }
  .case-card.reverse { direction: ltr; }
}

@media (max-width: 600px) {
  .work-grid  { grid-template-columns: 1fr; }
  .book-wrap  { padding: 48px 24px; }
  .hero-float-badge:nth-child(4) { display: none; }
}

@media (prefers-reduced-motion: reduce) {
  *, *::before, *::after {
    animation-duration: .001ms !important;
    transition-duration: .001ms !important;
  }
  .reveal, .stagger > * { opacity: 1; transform: none; }
  .hero-float-badge { animation: none; }
}
\end{lstlisting}

\begin{jargon}{Media query}
A block of CSS applying only when a condition about the display is true.
\code{@media (max-width: 960px)} means ``apply everything here only when the window
is 960 pixels wide or narrower''. Because it comes later in the file, its rules
override the equivalents above.
\end{jargon}

The 960-pixel block does the real work of making this design usable on a phone. The
hero collapses from two columns to one (line 540) and stops demanding the full
window height (line 541). The photo column becomes \code{55vw} tall with a
280-pixel floor (lines 545 and 546), so the photograph stays a sensible shape
instead of filling a tall narrow screen.

Lines 553 to 557 fold the metrics band from four columns to two and rebuild its
dividing lines: right border removed from everything, bottom border added to
everything, right border restored on odd items so there is still a line down the
middle, bottom border removed from items 3 and 4 so the block does not end in a
stray line. That is four rules to reposition two lines, and it is the price of the
border technique.

Line 558 collapses the case cards to a single column, and line 559 is the third dead
rule.

The 600-pixel block collapses the work grid to one column, tightens the booking
panel, and hides the third floating badge, which is the one positioned at the top
of the photo column where a phone has least room.

The final block, lines 568 to 575, is the best thing in the file.

\begin{keyidea}{What \code{prefers-reduced-motion} is, and why this block is better
than its sibling's}
Operating systems let people ask for reduced motion, usually because animation
causes nausea, dizziness or migraine. The browser passes that preference to the
page and this block responds.

It does three things. It crushes every animation and transition to a thousandth of
a millisecond, effectively instant rather than disabled, a deliberate choice because
setting them to zero can stop some browsers firing the events that mark an animation
as finished. Line 573 then forces every revealed and staggered element to full
opacity, because those start invisible and would otherwise never appear once their
transitions were neutered. That second line is the detail most implementations get
wrong.

Line 574 is the addition this file makes and the sibling does not:
\code{.hero-float-badge \{ animation: none; \}} stops the three badges drifting
outright. The blanket duration rule above would already have all but stopped them,
so this is belt and braces on the one continuous, looping animation on the page,
which is exactly the kind a motion-sensitive visitor would notice most.
\end{keyidea}

\begin{honest}{The reduced-motion block cannot reach the counting numbers}
This block governs CSS only. The four metric numbers are animated by JavaScript,
counting up over 2.2 seconds, and \file{main.js} never checks this preference. A
visitor who has asked for reduced motion still gets four numbers whirring upward.

The fix is four lines in \file{main.js}, given in section 7.3. This is the clearest
defect in the project after the placeholder: not because it is severe, but because
the intent to handle it correctly is demonstrably here, and expressed more carefully
than in the sibling, and still was not carried across to the other file.
\end{honest}

\subsubsection*{What it connects to}

\textbf{Named by:} \file{new-design-reference/index.html}, line 16, and nothing
else.

\textbf{Names:} no other file. It refers to no images and imports nothing. It names
font \emph{families}, which the HTML is responsible for downloading.

\textbf{Coupled to \file{index.html}} in three ways that no tool checks: by class
name throughout, by child position for the floating badges (lines 217 to 219), and
by the shared contract that JavaScript adds the class \code{visible}. One of those
couplings is already broken, at lines 436, 437 and 559.

\textbf{Coupled to \file{js/main.js}} through exactly two class names,
\code{scrolled} and \code{visible}, both of which the JavaScript adds and this file
defines.

\subsubsection*{Concepts introduced}

Custom properties, the cascade, flexbox, grid, pseudo-elements, pseudo-classes,
keyframe animations, transitions, media queries, \code{clamp}, \code{calc},
\code{aspect-ratio}, \code{backdrop-filter}, \code{background-clip: text}, stacking
with \code{z-index}, and \code{prefers-reduced-motion}. Each is defined at first use
above and collected in the glossary.

\subsubsection*{How it breaks}

\begin{itemize}
\item \textbf{Delete this file and the page still works}, in the sense that every
  word remains readable. It looks like an unstyled document, and the two-column
  hero becomes a stack.
\item \textbf{Rename a class here without renaming it in the HTML} and that element
  silently loses its styling. Lines 436, 437 and 559 are a live example.
\item \textbf{Insert an element before the placeholder in the HTML} and all three
  badges jump, because lines 217 to 219 position them by child number.
\item \textbf{Remove \code{pointer-events: none} from either fade} (lines 189, 198)
  and those invisible layers begin swallowing clicks.
\item \textbf{Remove \code{position: relative} on line 137} and the badges jump to
  positions measured against the whole window.
\item \textbf{A missing semicolon} makes the browser discard the rest of that
  declaration block. Unlike a program, CSS never reports an error: it drops what it
  cannot parse and carries on, which is why a stylesheet bug shows up as something
  looking wrong rather than as a message.
\end{itemize}

\subsubsection*{Honest findings for this file}

Five rules still using the agency blue after the rebrand (lines 65, 255, 256, 337,
435, 498); a sixth inherited literal in the badge background (line 205); four
tokens defined and never read; three dead rules for a class the HTML does not use
(lines 436, 437, 559); a doubled flex layout in the navigation; smooth scrolling
implemented twice; no focus states on either button; and a reduced-motion block
that is the best in either sibling and cannot reach the one animation living
outside CSS.
